% =============================================================================
% Lembar Kerja Fase 1: Fondasi Modern C++
% Universitas Bengkulu — Modern C++ Masterclass
% =============================================================================
\documentclass[12pt,a4paper]{article}
\usepackage[margin=2.5cm]{geometry}
\usepackage{amsmath,amssymb}
\usepackage{listings}
\usepackage{xcolor}
\usepackage{tcolorbox}
\usepackage{hyperref}
\usepackage{enumitem}

% Colors
\definecolor{UBgreen}{RGB}{0,102,51}
\definecolor{codebg}{RGB}{245,245,245}
\definecolor{keyword}{RGB}{0,0,180}
\definecolor{comment}{RGB}{0,130,0}
\definecolor{string}{RGB}{180,0,0}

\lstset{
  language=C++,
  backgroundcolor=\color{codebg},
  basicstyle=\ttfamily\small,
  keywordstyle=\color{keyword}\bfseries,
  commentstyle=\color{comment}\itshape,
  stringstyle=\color{string},
  frame=single,
  framerule=0.5pt,
  rulecolor=\color{UBgreen},
  breaklines=true,
  showstringspaces=false,
  numbers=left,
  numberstyle=\tiny\color{gray},
  numbersep=5pt
}

% Header
\newcommand{\worksheetheader}[2]{
  \noindent
  \begin{tabular}{@{}p{0.6\textwidth}p{0.4\textwidth}@{}}
    \textbf{\large Lembar Kerja: #1} & \textbf{Mata Kuliah:} Modern C++ untuk Power Systems \\
    \textbf{Nama:} \underline{\hspace{5cm}} & \textbf{NIM:} \underline{\hspace{3cm}} \\
    \textbf{Kelas:} \underline{\hspace{2cm}} & \textbf{Tanggal:} \underline{\hspace{2.5cm}}
  \end{tabular}
  \vspace{0.5cm}
  \hrule
  \vspace{0.5cm}
}

\begin{document}

\worksheetheader{Fase 1: consteval, string\_view, format, Structured Binding, Modules}

% =============================================================================
% BAGIAN A: KONSEP DASAR (C1 - Remember, C2 - Understand)
% =============================================================================
\section*{Bagian A: Konsep Dasar (C1--C2)}

\begin{tcolorbox}[colback=UBgreen!5, colframe=UBgreen, title=\textbf{A1. consteval vs constexpr (C1: Remember)}]
\textbf{Isilah tabel perbedaan berikut:}

\begin{center}
\begin{tabular}{|l|c|c|}
\hline
\textbf{Karakteristik} & \texttt{constexpr} & \texttt{consteval} \\
\hline
Bisa dievaluasi compile-time & Ya & Ya \\
\hline
\emph{Harus} dievaluasi compile-time & Tidak & \textbf{Ya} \\
\hline
Bisa dipanggil dengan argumen runtime & Ya & \textbf{Tidak} \\
\hline
Bisa digunakan di konteks non-constant & Ya & \textbf{Tidak} \\
\hline
Contoh use case & Fungsi utility umum & Validasi compile-time, lookup table \\
\hline
\end{tabular}
\end{center}

\textbf{Soal:} Mengapa \texttt{consteval} cocok untuk \texttt{fibonacci(10)} tapi tidak untuk \texttt{fibonacci(n)} di mana \texttt{n} input user?
\end{tcolorbox}

\vspace{0.5cm}

\begin{tcolorbox}[colback=UBgreen!5, colframe=UBgreen, title=\textbf{A2. std::string\_view (C2: Understand)}]
\textbf{Jelaskan dengan kalimat sendiri:}

\begin{enumerate}[label=\alph*)]
  \item Mengapa \texttt{std::string\_view} disebut \textit{zero-copy}?
  
  \vspace{2cm}
  
  \item Apa risiko \textit{dangling reference} saat menggunakan \texttt{string\_view}? Berikan contoh kode yang bermasalah.
  
  \vspace{3cm}
  
  \item Kapan sebaiknya gunakan \texttt{std::string} vs \texttt{std::string\_view}?
  
  \vspace{2cm}
\end{enumerate}
\end{tcolorbox}

% =============================================================================
% BAGIAN B: APLIKASI (C3 - Apply)
% =============================================================================
\section*{Bagian B: Aplikasi (C3 - Apply)}

\begin{tcolorbox}[colback=UBgreen!5, colframe=UBgreen, title=\textbf{B1. Parsing Log Grid (C3: Apply)}]
\textbf{Tulis fungsi \texttt{parse\_log\_line} yang menerima \texttt{std::string\_view} dan mengembalikan struct:}

\begin{lstlisting}
struct LogEntry {
    std::string_view timestamp;  // "2026-10-01T08:07:45"
    std::string_view severity;   // "CRITICAL"
    std::string_view message;    // "Transformer_T1 overloaded (115%)"
};

LogEntry parse_log_line(std::string_view line);
\end{lstlisting}

\textbf{Format input:} \texttt{2026-10-01T08:07:45 [CRITICAL] Transformer_T1 overloaded (115\%)}

\textbf{Petunjuk:} Gunakan \texttt{find()}, \texttt{substr()}, hindari allocation.

\vspace{4cm}
\textbf{Test cases:}
\begin{itemize}
  \item \texttt{"2026-10-01T08:00:00 [INFO] System started"}
  \item \texttt{"2026-10-01T08:05:12 [WARN] Voltage drop detected"}
  \item \texttt{"2026-10-01T08:07:45 [CRITICAL] Overload!"}
\end{itemize}
\end{tcolorbox}

\vspace{0.5cm}

\begin{tcolorbox}[colback=UBgreen!5, colframe=UBgreen, title=\textbf{B2. Structured Binding (C3: Apply)}]
\textbf{Lengkapi kode berikut untuk unpack struct \texttt{Measurement}:}

\begin{lstlisting}
struct Measurement {
    std::string sensor_name;
    double voltage;   // kV
    double current;   // A
};

int main() {
    Measurement m{"SUTET_150kV", 150.0, 200.0};
    
    // TODO: Structured binding di sini
    // auto [_____, _____, _____] = m;
    
    double apparent_power = _____ * _____; // VA
    std::cout << "Power: " << apparent_power << " VA\n";
}
\end{lstlisting}
\end{tcolorbox}

\vspace{0.5cm}

\begin{tcolorbox}[colback=UBgreen!5, colframe=UBgreen, title=\textbf{B3. std::format (C3: Apply)}]
\textbf{Gunakan \texttt{std::format} untuk output berikut:}

\begin{enumerate}[label=\alph*)]
  \item Tegangan: \texttt{230.5 kV}, Arus: \texttt{12.8 A} $\rightarrow$ \texttt{"V=230.5kV, I=12.8A"}
  \item Koordinat: \texttt{lat=-0.7893}, \texttt{lon=100.3542} $\rightarrow$ \texttt{"(-0.7893, 100.3542)"}
  \item Status: \texttt{true} $\rightarrow$ \texttt{"ONLINE"}, \texttt{false} $\rightarrow$ \texttt{"OFFLINE"}
\end{enumerate}

\begin{lstlisting}
// TODO: Lengkapi
std::cout << std::format("_______", _______);
\end{lstlisting}
\end{tcolorbox}

% =============================================================================
% BAGIAN C: ANALISIS (C4 - Analyze)
% =============================================================================
\section*{Bagian C: Analisis (C4 - Analyze)}

\begin{tcolorbox}[colback=UBgreen!5, colframe=UBgreen, title=\textbf{C1. Analisis Performa: string vs string\_view (C4: Analyze)}]
\textbf{Bandingkan dua implementasi parsing log 1 juta baris:}

\begin{lstlisting}
// Versi 1: std::string
std::string line;
while (std::getline(file, line)) {
    auto pos1 = line.find('[');
    auto pos2 = line.find(']');
    std::string severity = line.substr(pos1+1, pos2-pos1-1); // ALLOKASI!
    process(severity);
}

// Versi 2: std::string_view
std::string line;
while (std::getline(file, line)) {
    std::string_view sv = line;
    auto pos1 = sv.find('[');
    auto pos2 = sv.find(']');
    std::string_view severity = sv.substr(pos1+1, pos2-pos1-1); // NO ALLOKASI
    process(severity);
}
\end{lstlisting}

\textbf{Soal:}
\begin{enumerate}[label=\alph*)]
  \item Estimasi alokasi memori Versi 1 untuk 1 juta baris (rata-rata severity 7 char).
  
  \vspace{2cm}
  
  \item Mengapa Versi 2 masih butuh \texttt{std::string line} di luar loop?
  
  \vspace{2cm}
  
  \item Apa yang terjadi jika \texttt{process()} menyimpan \texttt{string\_view} ke vektor global?
  
  \vspace{2cm}
\end{enumerate}
\end{tcolorbox}

% =============================================================================
% BAGIAN D: EVALUASI (C5 - Evaluate)
% =============================================================================
\section*{Bagian D: Evaluasi (C5 - Evaluate)}

\begin{tcolorbox}[colback=UBgreen!5, colframe=UBgreen, title=\textbf{D1. Evaluasi Desain: consteval untuk Power Flow (C5: Evaluate)}]
\textbf{Scenario:} Anda ingin menghitung impedansi saluran standar IEC 60909 pada compile-time untuk konfigurasi grid yang \textbf{fixed} (diketahui saat compile).

\begin{itemize}
  \item Argumen \textbf{UNTUK} menggunakan \texttt{consteval}:
  
  \vspace{2cm}
  
  \item Argumen \textbf{TOLAK} menggunakan \texttt{consteval}:
  
  \vspace{2cm}
  
  \item \textbf{Keputusan Anda:} \hspace{2cm} \texttt{consteval} / \texttt{constexpr} / \texttt{const} / \texttt{runtime}
  \item \textbf{Alasan:}
  
  \vspace{2cm}
\end{itemize}
\end{tcolorbox}

\vspace{0.5cm}

\begin{tcolorbox}[colback=UBgreen!5, colframe=UBgreen, title=\textbf{D2. Evaluasi: C++20 Modules vs Header (C5: Evaluate)}]
\textbf{Bandingkan untuk proyek tim 5 orang, 50 file source:}

\begin{center}
\begin{tabular}{|l|c|c|}
\hline
\textbf{Aspek} & \textbf{Header (.hpp)} & \textbf{Modules (.cppm)} \\
\hline
Kompilasi bersyarat & Lambat (include berulang) & Cepat (BMI cached) \\
\hline
Enkapsulasi & Lemah (semua \texttt{public}) & Kuat (\texttt{export} control) \\
\hline
ODR violations & Rentan & Terhindar \\
\hline
Tooling support & Universal & Clang/MSVC (GCC terbatas) \\
\hline
Migrasi existing code & Mudah & Butuh refactor \\
\hline
\end{tabular}
\end{center}

\textbf{Keputusan untuk proyek baru 2025:} \hspace{3cm} Modules / Header
\end{tcolorbox}

% =============================================================================
% BAGIAN E: CREATIVE (C6 - Create)
% =============================================================================
\section*{Bagian E: Kreatif (C6 - Create)}

\begin{tcolorbox}[colback=UBgold!10, colframe=UBgold, title=\textbf{E1. Project Mini: Compile-Time Impedance Table (C6: Create)}]
\textbf{Buat program yang:}

\begin{enumerate}
  \item \texttt{consteval} generate \texttt{std::array<double, N>} impedansi saluran standar (misal: 10 tipe kabel).
  \item \texttt{consteval} validasi: impedansi $> 0$, tidak NaN.
  \item \texttt{main()} print table dengan \texttt{std::format} rapi.
  \item Bonus: \texttt{consteval} hitung \texttt{ampacity} berdasarkan temperatur ambient.
\end{enumerate}

\textbf{Template kode:}

\begin{lstlisting}
#include <array>
#include <format>
#include <iostream>

// TODO: Define CableSpec struct

// TODO: consteval array<CableSpec, 10> generate_standard_cables()

// TODO: consteval bool validate_cables(const auto& cables)

int main() {
    // TODO: Print table using std::format
    // No | Name        | R (ohm/km) | X (ohm/km) | Ampacity (A)
}
\end{lstlisting}

\vspace{1cm}
\textbf{Kriteria penilaian:}
\begin{itemize}
  \item[\square] Kompilasi berhasil (consteval dievaluasi compile-time)
  \item[\square] Validasi compile-time gagal jika data invalid
  \item[\square] Output tabel rapi dengan \texttt{std::format}
  \item[\square] Bonus: Ampacity calculation benar
\end{itemize}
\end{tcolorbox}

\vspace{0.5cm}

\begin{tcolorbox}[colback=UBgold!10, colframe=UBgold, title=\textbf{E2. Creative: Log Parser Module (C6: Create)}]
\textbf{Buat C++20 Module \texttt{grid\_parser} yang:}

\begin{enumerate}
  \item \texttt{export enum class Severity \{ INFO, WARN, CRITICAL, UNKNOWN \};}
  \item \texttt{export struct LogEntry \{ timestamp, severity, message \};}
  \item \texttt{export LogEntry parse\_log\_line(std::string\_view);}
  \item \texttt{export std::vector<LogEntry> parse\_file(std::string\_view path);}
  \item Handle edge cases: empty line, malformed, missing severity.
\end{enumerate}

\textbf{Test dengan \texttt{data/sample\_grid\_logs.txt}:}

\begin{lstlisting}
2026-10-01T08:00:00 [INFO] Substation_A initialization complete.
2026-10-01T08:05:12 [WARN] Bus_12 voltage dropping (0.94 pu).
2026-10-01T08:07:45 [CRITICAL] Transformer_T1 overloaded (115%).
...
\end{lstlisting}
\end{tcolorbox}

% =============================================================================
% REFLEKSI
% =============================================================================
\section*{Refleksi Diri}

\begin{tcolorbox}[colback=blue!5, colframe=blue, title=\textbf{Refleksi Fase 1}]
\begin{enumerate}[label=\arabic*.]
  \item Konsep mana yang paling sulit dipahami? \underline{\hspace{8cm}}
  
  \item Konsep mana yang paling berguna untuk proyek power systems Anda? \underline{\hspace{8cm}}
  
  \item Apa yang akan Anda terapkan ke proyek nyata minggu ini? \underline{\hspace{8cm}}
  
  \item Skor pemahaman (1-10): \underline{\hspace{1cm}}
\end{enumerate}
\end{tcolorbox}

\vspace{1cm}
\begin{center}
\textbf{\large --- Selamat Mengerjakan ---}
\end{center}

\end{document}